% =========================================================================
% PLANTILLA MAESTRA: CURRÍCULUM VITAE MODERNO (2 COLUMNAS / FOTO + QR)
% =========================================================================
% Características:
% - Maquetación asimétrica a dos columnas balanceadas con `paracol`.
% - Columna lateral con foto enmarcada en TikZ, idiomas, contacto y código QR.
% - Tipografía limpia Sans-Serif (Helvetica / Arial).
% - Íconos vectoriales y banderas diseñados nativamente en TikZ (sin dependencias).
% - Ideal para puestos que valoran fuerte presencia visual o perfil comercial/tech.
% =========================================================================

\documentclass[10pt,a4paper]{article}

% --- PAQUETES DE CONFIGURACIÓN Y TIPOGRAFÍA ---
\usepackage[utf8]{inputenc}
\usepackage[T1]{fontenc}
\usepackage[spanish,es-nodecimaldot]{babel}
\usepackage{helvet}
\renewcommand{\familydefault}{\sfdefault}
\usepackage[a4paper,top=1.15cm,bottom=1.15cm,left=1.2cm,right=1.2cm]{geometry}
\usepackage{paracol}
\usepackage{amssymb}
\usepackage{tikz}
\usetikzlibrary{calc}
\usepackage{graphicx}
\usepackage{xcolor}
\usepackage{enumitem}
\usepackage{tcolorbox}
\tcbuselibrary{skins}
\usepackage{microtype}
\usepackage{hyperref}

\pagestyle{empty}

% --- PALETA DE COLORES PROFESIONAL ---
\definecolor{primary}{RGB}{26, 54, 93}       % Azul marino corporativo (#1A365D)
\definecolor{secondary}{RGB}{43, 108, 176}   % Azul pizarra medio (#2B6CB0)
\definecolor{darktext}{RGB}{30, 41, 59}      % Carbón oscuro para lectura (#1E293B)
\definecolor{muted}{RGB}{100, 116, 139}      % Gris medio para fechas y metadatos (#64748B)
\definecolor{lightbg}{RGB}{248, 250, 252}    % Fondo claro sutil (#F8FAFC)
\definecolor{bordercolor}{RGB}{226, 232, 240}% Gris suave para divisores (#E2E8F0)

\hypersetup{
    colorlinks=true,
    linkcolor=primary,
    urlcolor=secondary,
    pdfauthor={[NOMBRE CANDIDATO]},
    pdftitle={CV - [NOMBRE CANDIDATO]},
    pdfsubject={Curriculum Vitae},
    pdfcreator={LaTeX - pdflatex}
}

% --- MACROS Y FORMATO DE SECCIONES ---
\newcommand{\cvsection}[1]{%
  \vspace{5pt}%
  \noindent{\color{primary}\bfseries\large\MakeUppercase{#1}}\par%
  \vspace{2pt}%
  \noindent{\color{secondary}\rule{\linewidth}{1.2pt}}\par%
  \vspace{4pt}%
}

\newcommand{\cvsidesection}[1]{%
  \vspace{6pt}%
  \noindent{\color{primary}\bfseries\normalsize\MakeUppercase{#1}}\par%
  \vspace{2pt}%
  \noindent{\color{secondary}\rule{\linewidth}{1.0pt}}\par%
  \vspace{4pt}%
}

% --- ÍCONOS VECTORIALES TIKZ (NATIVOS Y ESCALABLES) ---
\newcommand{\iconEmail}{%
  \begin{tikzpicture}[x=0.07cm, y=0.07cm, baseline=-0.6ex]
    \draw[fill=primary!12, draw=primary, line width=0.5pt, rounded corners=0.4pt] (0,0) rectangle (11,8);
    \draw[primary, line width=0.5pt] (0,8) -- (5.5,3.8) -- (11,8);
  \end{tikzpicture}%
}

\newcommand{\iconPhone}{%
  \begin{tikzpicture}[x=0.07cm, y=0.07cm, baseline=-0.6ex]
    \draw[fill=primary!12, draw=primary, line width=0.5pt, rounded corners=0.8pt] (0,0) rectangle (6.5,11);
    \draw[primary, line width=0.5pt] (2,9.6) -- (4.5,9.6);
    \fill[primary] (3.25,1.5) circle (0.6pt);
    \draw[fill=white, draw=primary, line width=0.3pt] (0.9,2.8) rectangle (5.6,8.6);
  \end{tikzpicture}%
}

\newcommand{\iconLocation}{%
  \begin{tikzpicture}[x=0.07cm, y=0.07cm, baseline=-0.8ex]
    \fill[primary] (3.5,8) circle (3);
    \fill[white] (3.5,8) circle (1.1);
    \fill[primary] (0.8,6.8) -- (3.5,0.5) -- (6.2,6.8) -- cycle;
  \end{tikzpicture}%
}

\newcommand{\iconLink}{%
  \begin{tikzpicture}[x=0.07cm, y=0.07cm, baseline=-0.6ex]
    \draw[primary, line width=0.6pt, rounded corners=0.8pt] (0,2.5) rectangle (7,7);
    \draw[primary, line width=0.6pt] (7,4.75) -- (10.5,4.75);
    \draw[primary, line width=0.6pt, rounded corners=0.8pt] (3.5,0) rectangle (10.5,4.5);
  \end{tikzpicture}%
}

\newcommand{\flagUS}{%
  \begin{tikzpicture}[x=0.08cm, y=0.08cm, baseline=-0.5ex]
    \draw[draw=gray!30, line width=0.3pt] (0,0) rectangle (11,7.5);
    \clip (0,0) rectangle (11,7.5);
    \foreach \i in {0,2,4,6} {
      \fill[red!80!black] (0,\i) rectangle (11,\i+0.95);
    }
    \foreach \i in {1,3,5} {
      \fill[white] (0,\i) rectangle (11,\i+0.95);
    }
    \fill[blue!70!black] (0,3.8) rectangle (5.5,7.5);
  \end{tikzpicture}%
}

\newcommand{\flagES}{%
  \begin{tikzpicture}[x=0.08cm, y=0.08cm, baseline=-0.5ex]
    \draw[draw=gray!30, line width=0.3pt] (0,0) rectangle (11,7.5);
    \fill[red!80!black] (0,0) rectangle (11,7.5);
    \fill[yellow!90!black] (0,2) rectangle (11,5.5);
  \end{tikzpicture}%
}

% --- MACROS DE ENTRADAS ---
\newcommand{\cvexperience}[5]{%
  \noindent{\color{primary}\textbf{\normalsize #1}}\hfill{\footnotesize\color{muted}\textbf{#2}}\par
  \noindent{\small\color{secondary}\textbf{#3}}\par
  \if\relax\detokenize{#4}\relax\else
    \vspace{1.5pt}\noindent{\footnotesize\itshape\color{darktext}#4}\par
  \fi
  \vspace{2pt}
  #5
  \vspace{4.5pt}
}

\begin{document}
\color{darktext}

% ==============================================================================
% PÁGINA 1
% ==============================================================================

\setcolumnwidth{5.6cm, 12.2cm}
\setlength{\columnsep}{0.8cm}
\setlength{\columnseprule}{0.5pt}
\colseprulecolor{bordercolor}

\begin{paracol}{2}

% --- PÁGINA 1: COLUMNA IZQUIERDA ---

% Fotografía de perfil enmarcada (reemplazar foto_perfil.png o comentar bloque)
\begin{center}
  \begin{tikzpicture}
    \clip[rounded corners=6pt] (0,0) rectangle (4.2cm, 5.6cm);
    \node[inner sep=0pt, anchor=south west] at (0,0) {%
      \IfFileExists{foto_perfil.png}{\includegraphics[width=4.2cm, height=5.6cm]{foto_perfil.png}}{\rule{4.2cm}{5.6cm}}%
    };
  \end{tikzpicture}
\end{center}
\vspace{-4pt}

% Formación Académica
\cvsidesection{Formación}

\noindent\textbf{\small [Título / Grado Académico]}\par
\noindent{\footnotesize\color{secondary}[Universidad o Facultad]}\par
\noindent{\scriptsize\color{muted}[Año / Estado]}\vspace{4pt}

\noindent\textbf{\small [Título Secundario o Terciario]}\par
\noindent{\footnotesize\color{secondary}[Institución Educativa]}\par
\noindent{\scriptsize\color{muted}[Año]}\vspace{4pt}

\noindent\textbf{\small [Certificación Clave]}\par
\noindent{\footnotesize\color{secondary}[Entidad Emisora]}\par
\noindent{\scriptsize\color{muted}[Año]}\vspace{4pt}

% Idiomas
\cvsidesection{Idiomas}
\begin{itemize}[leftmargin=14pt, itemsep=2.5pt, topsep=1pt]
  \item[\flagUS] \textbf{\small Inglés:} [Nivel / B2]
  \item[\flagES] \textbf{\small Español:} Nativo
\end{itemize}
\vspace{2pt}

% Contacto
\cvsidesection{Contacto}
\begin{itemize}[leftmargin=14pt, itemsep=3.5pt, topsep=1pt]
  \item[\iconEmail] {\scriptsize\href{mailto:correo@ejemplo.com}{correo@ejemplo.com}}
  \item[\iconPhone] {\scriptsize\href{tel:+549000000000}{(+54) 000-000000}}
  \item[\iconLocation] {\scriptsize [Ciudad, País]}
  \item[\iconLink] {\scriptsize\href{https://linktr.ee/usuario}{\textbf{Perfil / Portafolio}}}
\end{itemize}

\vspace{5pt}
% Código QR
\begin{center}
  \begin{tikzpicture}
    \node[draw=bordercolor, fill=white, rounded corners=4pt, inner sep=3pt] {
      \IfFileExists{qr_linktree.png}{\includegraphics[width=2.2cm]{qr_linktree.png}}{\rule{2.2cm}{2.2cm}}
    };
  \end{tikzpicture}\par
  \vspace{2pt}
  {\scriptsize\color{muted}[Enlace o Portafolio]}
\end{center}

\switchcolumn

% --- PÁGINA 1: COLUMNA DERECHA ---

% Cabecera Principal
{\fontsize{21}{25}\selectfont\color{primary}\textbf{[NOMBRE COMPLETO]}}\par
\vspace{3pt}
{\fontsize{10.5}{13}\selectfont\color{secondary}\textbf{[TÍTULO PROFESIONAL O ESPECIALIDAD]}}\par
\vspace{5pt}

% Acerca de Mí
\cvsection{Acerca de Mí}
\begin{tcolorbox}[
  width=\linewidth,
  colback=lightbg,
  colframe=secondary!40,
  boxrule=0pt,
  leftrule=2.5pt,
  arc=2pt,
  top=4pt,bottom=4pt,left=6pt,right=6pt
]
  \small\itshape\color{darktext}
  ``[Párrafo de presentación profesional destacando formación, habilidades técnicas diferenciales, enfoque metodológico y objetivos laborales clave].''
\end{tcolorbox}
\vspace{2pt}

% Experiencia Profesional
\cvsection{Experiencia Profesional}

\cvexperience{[Puesto / Rol]}{[Periodo]}{[Empresa / Institución]}{[Resumen ejecutivo de la función y logros alcanzados]}{%
  \begin{itemize}[leftmargin=11pt, itemsep=1.5pt, topsep=1pt]
    \item [Responsabilidad o hito técnico clave cuantificado].
    \item [Herramientas, tecnologías y metodologías aplicadas].
    \item [Resultados de negocio o eficiencia operativa generada].
  \end{itemize}
}

\cvexperience{[Puesto Anterior]}{[Periodo]}{[Empresa]}{[Descripción breve de funciones]}{%
  \begin{itemize}[leftmargin=11pt, itemsep=1.5pt, topsep=1pt]
    \item [Tarea o proyecto desempeñado].
    \item [Resolución de problemas técnicos].
  \end{itemize}
}

\end{paracol}

% ==============================================================================
% PÁGINA 2
% ==============================================================================
\newpage

\begin{paracol}{2}

% --- PÁGINA 2: COLUMNA IZQUIERDA ---

\cvsidesection{Habilidades}

\begin{itemize}[leftmargin=11pt, itemsep=4pt, topsep=2pt, label=\textcolor{secondary}{\scriptsize$\blacktriangleright$}]
  \item \textbf{Habilidad Blanda 1:} [Descripción breve].
  \item \textbf{Habilidad Blanda 2:} [Descripción breve].
  \item \textbf{Habilidad Blanda 3:} [Descripción breve].
  \item \textbf{Habilidad Blanda 4:} [Descripción breve].
  \item \textbf{Habilidad Blanda 5:} [Descripción breve].
  \item \textbf{Habilidad Blanda 6:} [Descripción breve].
\end{itemize}

\vspace{18pt}
\cvsidesection{Contacto Rápido}

\begin{itemize}[leftmargin=14pt, itemsep=4pt, topsep=1pt]
  \item[\iconEmail] {\scriptsize\href{mailto:correo@ejemplo.com}{correo@ejemplo.com}}
  \item[\iconPhone] {\scriptsize\href{tel:+549000000000}{(+54) 000-000000}}
  \item[\iconLocation] {\scriptsize [Ciudad, País]}
  \item[\iconLink] {\scriptsize\href{https://linktr.ee/usuario}{\textbf{Portafolio}}}
\end{itemize}

\switchcolumn

% --- PÁGINA 2: COLUMNA DERECHA ---

% Proyectos Destacados
\cvsection{Proyectos Académicos y Personales}

\noindent{\color{primary}\textbf{\normalsize [Nombre del Proyecto 1]}}\hfill{\footnotesize\color{muted}\textbf{[Año]}}\par
\vspace{1.5pt}
\begin{itemize}[leftmargin=11pt, itemsep=2pt, topsep=1pt]
  \item \textbf{Objetivo y Diseño:} [Descripción del alcance y solución técnica].
  \item \textbf{Tecnologías:} [Lenguajes, frameworks y bases de datos empleadas].
\end{itemize}
\vspace{5pt}

% Conocimientos Técnicos
\cvsection{Conocimientos Técnicos}

\begin{tcolorbox}[
  width=\linewidth,
  colback=white,
  colframe=bordercolor,
  boxrule=0.6pt,
  arc=3pt,
  top=6pt,bottom=6pt,left=8pt,right=8pt
]
\small
\noindent\textbf{\color{primary}Sistemas Operativos:} [Windows, Linux, macOS].\par\vspace{3.5pt}
\noindent\textbf{\color{primary}Lenguajes de Programación:} [Python, Java, C++, etc.].\par\vspace{3.5pt}
\noindent\textbf{\color{primary}Bases de Datos:} [PostgreSQL, MySQL, SQLite].\par\vspace{3.5pt}
\noindent\textbf{\color{primary}Herramientas y DevOps:} [Docker, Git, Bash, etc.].
\end{tcolorbox}

\end{paracol}

\end{document}
